\section{Relation to gyroscopic memory}
\subsection{Polarization area of a weak TT wave}
On the same local radiation-zone patch with the adapted screen fixed in the preceding section, write the leading radiative data at a large reference Bondi radius $r$ as a weak TT strain
\begin{equation}
h^{TT}(u,r;\eps)=\eps\begin{pmatrix}p(u)&q(u)\\q(u)&-p(u)\end{pmatrix}+O(\eps^2).
\label{eq:tt}
\end{equation}
In the radiation zone, $p,q=O(r^{-1})$. We work at fixed $r$ and suppress the $r$ argument below. Stationarity before and after the burst implies $\dot p(u_\pm)=\dot q(u_\pm)=0$. If
\[
X(\eps)=X^{(0)}+\eps X^{(1)}+\eps^2X^{(2)}+O(\eps^3),
\]
we denote the coefficient of $\eps^2$ by $X^{(2)}$. Thus $\Delta\Psi_{\rm H}^{(2)}$ is neither the full finite-amplitude memory nor the full order-$\eps^2$ contribution with the explicit factor $\eps^2$ included; it is the quadratic perturbative coefficient.

With metric signature $(-+++)$, the convention $T_{AB}=-R(e_A,k,e_B,k)$ from Eq.~\eqref{eq:tidal}, and our outgoing/retarded-time convention, a radiation-crossing null line normalized by $du/d\lambda=1$ and equipped with a screen adapted to the Bondi/TT polarization dyad has leading type-N radiative curvature
\[
R(e_A,k,e_B,k)=-\frac12\ddot h^{TT}_{AB}
+\text{(subleading in $1/r$ and nonlinear amplitude)}.
\]
Hence the leading tidal profile is $T=\frac12\ddot h^{TT}$, so the same pair $(p,q)$ can be inserted into both the Jacobi and Bondi/TT descriptions. The required statement is this fixed local identification; we do not assert that arbitrary null directions at finite $r$ give the same $T$. The overall sign of the TT/Bondi relation is convention dependent, but a simultaneous sign flip of $(p,q)$ leaves the quadratic area below unchanged. We fix the sign convention in Eq.~\eqref{eq:shear-strain} and use it consistently.

The quadratic Jacobi coefficient is the signed polarization-plane area obtained by closing the trajectory $(p,q)$ with a straight chord,
\begin{equation}
\mathcal A_{\rm pol}[p,q]
=\frac12\int_{u_-}^{u_+}(p\dot q-q\dot p)\,du
+\frac12[p_fq_i-p_iq_f].
\label{eq:area}
\end{equation}
More precisely,
\begin{equation}
\omega_\gamma(\eps)=\eps^2\omega_\gamma^{(2)}+O(\eps^3),
\qquad
\omega_\gamma^{(2)}=\mathcal A_{\rm pol}[p,q].
\label{eq:omega-area}
\end{equation}
In the initially shear-free reference frame, $p_i=q_i=0$, so the chord term vanishes, while the final shear may remain nonzero.

\subsection{Factor-of-two relation to the Bondi hard functional}
On the local screen write
\[
C=\begin{pmatrix}X&Y\\Y&-X\end{pmatrix},\qquad
\widetilde C=\begin{pmatrix}-Y&X\\X&Y\end{pmatrix}.
\]
We use $\epsilon_{12}=+1$ and the same index ordering as Seraj--Oblak and Faye--Seraj,
$\widetilde C_{AB}=\epsilon_{CA}C^C{}_B$. Reversing the index order in the dual convention would flip the overall sign, but that is not the convention used here. Therefore
\begin{equation}
N_{AB}\widetilde C^{AB}=2(X\dot Y-Y\dot X).
\label{eq:stf-algebra}
\end{equation}
Following the polarization convention of Faye--Seraj, we take in the radiation zone
\begin{equation}
C_{AB}=-r h^{TT}_{AB}+O(r^0).
\label{eq:shear-strain}
\end{equation}
Eqs.~\eqref{eq:hard-memory}, \eqref{eq:tt}, and \eqref{eq:stf-algebra} give
\begin{equation}
\Delta\Psi_{\rm H}(\eps)
=\eps^2\Delta\Psi_{\rm H}^{(2)}+O(\eps^3),
\qquad
\Delta\Psi_{\rm H}^{(2)}
=-\frac14\int(p\dot q-q\dot p)\,du.
\label{eq:gyro-area}
\end{equation}
Under the initially shear-free condition, Eq.~\eqref{eq:omega-area} becomes
\begin{equation}
\omega_\gamma^{(2)}=\frac12\int(p\dot q-q\dot p)\,du.
\label{eq:jacobi-area-eps}
\end{equation}
Hence:

\begin{theorem}[Factor relation between the two quadratic functionals of the same radiative data]
\label{thm:bridge}
In the standard initially shear-free Bondi frame, at weak amplitude and leading radiation-zone order, and with the oriented-screen conventions above,
\begin{equation}
\boxed{\omega_\gamma^{(2)}=-2\Delta\Psi_{\rm H}^{(2)}},\qquad
\boxed{|\omega_\gamma^{(2)}|=2|\Delta\Psi_{\rm H}^{(2)}|}.
\label{eq:bridge}
\end{equation}
\end{theorem}

The quantity $\omega_\gamma^{(2)}$ is a Jacobi scattering observable on a null geodesic, whereas $\Delta\Psi_{\rm H}^{(2)}$ is an orientation angle measured on a timelike gyroscope worldline; they are not identified as the same observable. The equality states that, under the local radiation-zone patch, radiation-crossing null direction, and adapted-screen identification fixed above, the two known quadratic functionals evaluated on the same local TT/Bondi radiative data $(p,q)$ differ by the factor shown. A different null direction or screen identification defines a new auxiliary-probe problem for which the relation must be rechecked. Reversing screen orientation flips both pseudoscalars simultaneously. We do not claim a theorem that controls finite $r$ and finite amplitude simultaneously; the safe coefficient statement is the iterated limit
\begin{equation}
\lim_{r\to\infty}r^2\lim_{\eps\to0}\eps^{-2}
\left(\omega_\gamma(\eps,r)+2\Delta\Psi_{\rm H}(\eps,r)\right)=0.
\label{eq:iterated-limit}
\end{equation}

\section{Stationary physical minimum-action theorem}
Endpoint stationarity determines the optimal coefficient. On the normalized interval, write the first-order tidal profile as
\[
A^{(1)}(t)=\alpha(t)\sigma_3+\beta(t)\sigma_1,\qquad
\alpha=\frac12p'',\quad \beta=\frac12q''.
\]
Then $A(t;\eps)=\eps A^{(1)}(t)+O(\eps^2)$. At fixed $r$, $p,q=O(r^{-1})$, so $\alpha,\beta=O(r^{-1})$ and both $Q_\gamma^{(2)}$ and $\Delta\Psi_{\rm H}^{(2)}$ scale as $O(r^{-2})$. Therefore no additional factor of $r$ remains in the ratio or in the optimal coefficient. The Bel--Robinson curvature action has expansion
\begin{equation}
Q_\gamma(\eps)=\eps^2Q_\gamma^{(2)}+O(\eps^3),\qquad
Q_\gamma^{(2)}=\|\alpha\|_2^2+\|\beta\|_2^2.
\label{eq:Q-leading}
\end{equation}
while $p'(0)=p'(1)=q'(0)=q'(1)=0$ implies
\begin{equation}
\int_0^1\alpha(t)\,dt=\int_0^1\beta(t)\,dt=0.
\label{eq:mean-zero}
\end{equation}
Introduce
\begin{equation}
\mathcal U_0:=\left\{f\in L^2(0,1):\int_0^1 f(t)\,dt=0\right\},\qquad
P_0f=f-\int_0^1f(s)\,ds,
\label{eq:U0}
\end{equation}
and define
\begin{equation}
K_{\rm stat}:=P_0K_0P_0,
\qquad
s_{\rm stat}:=\|K_{\rm stat}\|_{\mathcal U_0\to\mathcal U_0}.
\label{eq:Kstat}
\end{equation}
The operator $K_{\rm stat}$ is compact and skew-adjoint on $\mathcal U_0$.

Define the stationary exact-Jacobi value function by
\begin{equation}
\kappa_{\rm stat}(\omega):=
\inf\{Q_\gamma[A]:\omega_\gamma[A]=\omega,\ \alpha,\beta\in\mathcal U_0\}.
\label{eq:kappa-stat}
\end{equation}

\begin{theorem}[Sharp law for the stationary Jacobi channel]
\label{thm:stationary-sharp}
\begin{equation}
\boxed{\kappa_{\rm stat}(\omega)=\frac{|\omega|}{s_{\rm stat}}+O(\omega^2),\qquad \omega\to0.}
\label{eq:stationary-sharp}
\end{equation}
The coefficient $s_{\rm stat}^{-1}$ is sharp for the exact endpoint problem.
\end{theorem}
\begin{proof}
For $\alpha,\beta\in\mathcal U_0$,
\[
\langle\alpha,K_0\beta\rangle
=\langle\alpha,P_0K_0P_0\beta\rangle
=\langle\alpha,K_{\rm stat}\beta\rangle.
\]
Hence
\[
2|\langle\alpha,K_0\beta\rangle|
\le s_{\rm stat}(\|\alpha\|^2+\|\beta\|^2)
=s_{\rm stat}Q_\gamma.
\]
The remainder estimate from the unrestricted theorem is unchanged on the closed subspace, so the lower-bound argument for Theorem~\ref{thm:jacobi-sharp} applies with $s_0$ replaced by $s_{\rm stat}$.

For sharpness, choose a real unit eigenvector $\beta_*$ of the top eigenvalue $s_{\rm stat}^2$ of $-K_{\rm stat}^2$ on the real Hilbert space $\mathcal U_0$, and set
\[
\alpha_*:=K_{\rm stat}\beta_*/s_{\rm stat}.
\]
Then
\[
\|\alpha_*\|=\|\beta_*\|=1,\qquad
\langle\alpha_*,K_{\rm stat}\beta_*\rangle=s_{\rm stat}.
\]
With $A_*=\alpha_*\sigma_3+\beta_*\sigma_1$ and $A_\rho=\rho A_*$, the mean-zero constraints are preserved for every $\rho$, and
\[
Q_\gamma[A_\rho]=2\rho^2,\qquad
\omega_\gamma[A_\rho]=2s_{\rm stat}\rho^2+O(\rho^4).
\]
The same one-sided local inversion as in Theorem~\ref{thm:jacobi-sharp}, together with a sign flip of one polarization for negative targets, exactly recovers sufficiently small targets of either sign. The upper and lower asymptotics therefore agree, proving Eq.~\eqref{eq:stationary-sharp}.
\end{proof}

\begin{lemma}[Realization by stationary TT strain]
\label{lem:physical-realization}
For any $\alpha,\beta\in\mathcal U_0$, define
\begin{equation}
p(t)=2\int_0^t(t-s)\alpha(s)\,ds,\qquad
q(t)=2\int_0^t(t-s)\beta(s)\,ds.
\label{eq:pq-from-control}
\end{equation}
Then
\[
p(0)=q(0)=p'(0)=q'(0)=p'(1)=q'(1)=0,
\qquad p''=2\alpha,\ q''=2\beta.
\]
Thus the controls lift to first-order TT data that are initially shear free and stationary both before and after the burst.
\end{lemma}
\begin{proof}
The initial conditions follow directly from the definition. Since $p'(t)=2\int_0^t\alpha(s)ds$, one has $p'(1)=2\int_0^1\alpha=0$, and similarly for $q$. The final values $p(1)$ and $q(1)$ are generally nonzero, so a constant final shear is allowed.
\end{proof}

\begin{lemma}[The same sharp constant for smooth stationary bursts]
\label{lem:smooth-density}
Let
\[
\mathcal U_{0,c}^\infty:=\left\{f\in C_c^\infty(0,1):\int_0^1f(t)\,dt=0\right\}.
\]
Then $\mathcal U_{0,c}^\infty$ is dense in $\mathcal U_0$. Consequently, restricting the admissible class to smooth stationary TT bursts whose curvature vanishes near the endpoints does not change the sharp leading-order coefficient $s_{\rm stat}$.
\end{lemma}
\begin{proof}
For $f\in\mathcal U_0$, choose $g_n\in C_c^\infty(0,1)$ with $g_n\to f$ in $L^2$. Then $\int g_n\to0$. Fix $\chi\in C_c^\infty(0,1)$ with $\int\chi=1$ and define
\[
f_n:=g_n-\left(\int_0^1g_n\right)\chi.
\]
Then $f_n\in\mathcal U_{0,c}^\infty$ and $f_n\to f$ in $L^2$. Since $K_{\rm stat}$ is bounded, the bilinear form $\langle\alpha,K_{\rm stat}\beta\rangle$ is continuous on $L^2\times L^2$. Hence the top singular pair can be approximated arbitrarily well by smooth mean-zero functions. Each approximating pair vanishes near the endpoints in curvature, so the corresponding $p,q$ from Eq.~\eqref{eq:pq-from-control} extend smoothly as constants outside the active interval. The infimum over the smooth stationary-burst class therefore has the same sharp constant.
\end{proof}

\section{Main theorem: sharp lower bound on Bel--Robinson curvature action}
For a leading-order hard-memory target $\psi=\Delta\Psi_{\rm H}^{(2)}$, define the value function of the \emph{leading-order radiative-data problem} by
\begin{equation}
\kappa_{\rm H}^{\rm lead}(\psi):=
\inf\left\{Q_\gamma^{(2)}[A^{(1)}]:
\Delta\Psi_{\rm H}^{(2)}[A^{(1)}]=\psi,\ \alpha,\beta\in\mathcal U_0\right\}.
\label{eq:kappa-hard-lead}
\end{equation}
The minimization is over the weak-amplitude TT/Bondi tangent-data class above; no value function over the full space of finite-amplitude Einstein solutions is defined.

\begin{theorem}[Sharp minimum curvature action for a fixed radiation-crossing null window]
\label{thm:main-physical}
Assume the Bel--Robinson convention of Sec.~II, a fixed radiation-crossing null window $\gamma$ and adapted screen in a local radiation-zone patch with $L/r\to0$, the standard initially shear-free Bondi frame, weak amplitude and leading radiation-zone order, and stationary-in/stationary-out conditions, without imposing zero final shear. Then
\begin{equation}
\boxed{\kappa_{\rm H}^{\rm lead}(\psi)=\frac{2}{\|K_{\rm stat}\|}|\psi|.}
\label{eq:main-physical-value}
\end{equation}
This equality is exact for the leading-order perturbative radiative-data problem, and $2/\|K_{\rm stat}\|$ is the optimal coefficient.
\end{theorem}
\begin{proof}
The quadratic leading term in Eq.~\eqref{eq:endpoint-expansion}, combined with Theorem~\ref{thm:bridge}, gives in the initially shear-free stationary class
\[
\Delta\Psi_{\rm H}^{(2)}=-\langle\alpha,K_{\rm stat}\beta\rangle.
\]
Therefore
\[
2|\psi|=2|\langle\alpha,K_{\rm stat}\beta\rangle|
\le s_{\rm stat}(\|\alpha\|_2^2+\|\beta\|_2^2)
=s_{\rm stat}Q_\gamma^{(2)},
\]
so $Q_\gamma^{(2)}\ge2|\psi|/s_{\rm stat}$.

Conversely, use the real top singular pair $\alpha_*,\beta_*$ constructed in the proof of Theorem~\ref{thm:stationary-sharp}. Choose the sign of $\beta_*$ according to the sign of the target and set
\[
\alpha=r\alpha_*,\qquad \beta=\pm r\beta_*,\qquad
r^2=\frac{|\psi|}{s_{\rm stat}}.
\]
Then $\Delta\Psi_{\rm H}^{(2)}=\psi$ and
\[
Q_\gamma^{(2)}=2r^2=\frac{2}{s_{\rm stat}}|\psi|
\]
exactly. Lemma~\ref{lem:physical-realization} lifts the controls to initially shear-free, stationary-in/stationary-out TT data, and Lemma~\ref{lem:smooth-density} shows that the same infimum is obtained after restricting to smooth bursts whose curvature vanishes near the endpoints.
\end{proof}

Thus, for the fixed null window and within the stated leading-order class,
\begin{equation}
Q_\gamma^{(2)}\ge\frac{2}{\|K_{\rm stat}\|}|\Delta\Psi_{\rm H}^{(2)}|.
\label{eq:main-bound}
\end{equation}
Numerically,
\begin{equation}
\|K_{\rm stat}\|=0.0281140680170\ldots,
\qquad
\boxed{\frac{2}{\|K_{\rm stat}\|}=71.13876223\ldots.}
\label{eq:main-bound-num}
\end{equation}
These decimal values are validation only and do not define the theorem constant.

\begin{remark}[Relation to the earlier coefficient $21.96874242\ldots$]
The unrestricted $L^2$ coefficient $2/\|K_0\|=21.96874242\ldots$ remains a valid weaker lower bound on the stationary subclass, but it is not sharp there. Under the mean-zero constraints the optimization must be redone for $K_{\rm stat}=P_0K_0P_0$, yielding the optimal coefficient $2/\|K_{\rm stat}\|$ above.
\end{remark}